\scanpage{471}
\mspage{238}

not to the situation, perhaps submit it to Carlos, who is at ease both with discrete groups and their generators and relations, and with the theory of $\frac{1}{2}$ simples -- and topology ....

\scanpage{472}
\mspage{239}

\uline{33} Study of finite coverings -- relation between the $\mathcal N_{g,\nu},\Gamma_{g,\nu}$ for variable $g$.

Let $(\pi\deleted{,\unsure{a}})$ be an arithmetized profinite outer group with loops of type $(g,\nu)$, $\pi'$ an open \unsure{subgroup|subgroup} of $\pi$, whence an injective map
\[
\mathrm{Discr\acute{e}t}(\pi)\hookrightarrow
\mathrm{Discr\acute{e}t}(\pi'),
\]
we would like to prove that it is compatible with the equivalence relation \deleted{\unsure{of the}\ \illegible{1}}\inserted{\unsure{of arithmetization|of arithmetizability}}, in \deleted{\unsure{prove}} such a way that every arithmetization of $\pi$ gives one of $\pi'$, and likewise for the \deleted{\unsure{discre}} prediscretifications.

We would like to establish at the same time that the maps
\[
\begin{array}{lll}
 &P^+(\pi)\longrightarrow P^+(\pi')&\begin{gathered}\text{on the oriented}\\\text{prediscretifications}\end{gathered}\\[1em]
\text{and}&A(\pi)\unsure{\longrightarrow|-}A(\pi')&\text{on the arithmetizations}
\end{array}
\]
are injective. To prove these points, we are reduced to the case where $\pi'$ is characteristically invariant. (Cf. \S\,28, diagram 7, \unsure{ff}) Choosing

\scanpage{473}

a discretification $\pi_0$ of $\pi$, hence $\pi'_0$ of $\pi'$, we then find
\[
\unsure{\widehat{\widehat{\mathcal G}}}(\pi_0)\longrightarrow
\unsure{\widehat{\widehat{\mathcal G}}}(\pi'_0)
\]
and we have

\textcircled{1} The homomorphism $\unsure{\widehat{\widehat{\mathcal G}}}(\pi_0)\unsure{\longrightarrow}\unsure{\widehat{\widehat{\mathcal G}}}(\pi'_0)$ sends $\mathcal M(\pi_0)$ into $\mathcal M(\pi'_0)$

we will therefore have a diagram (a \unsure{variant|relative} of the one in \S\,28)
\diagram{\begin{tikzcd}[ampersand replacement=\&,column sep=large,row sep=large]
\unsure{(1)} \&
\widehat{\mathcal G}^{+}(\pi_0)
 \arrow[r,hook,no head] \arrow[d,hook] \&
\mathcal M(\pi_0) \arrow[r,hook,no head] \arrow[d,hook] \&
\widehat{\widehat{\mathcal G}}(\pi_0) \arrow[d,hook] \\
\& \widehat{\mathcal G}^{+}(\pi'_0)
 \arrow[r,hook,no head] \&
\mathcal M^{\unsure{\ast|\sharp}}(\pi'_0) \arrow[r,hook,no head] \&
\unsure{\widehat{\widehat{\mathcal G}}}(\pi'_0)
\end{tikzcd}}
which \deleted{means} implies that a discretification $\pi_1$ on $\pi$ which gives the same oriented prediscretification as $\pi_0$ (resp. the same arithmetization) defines a \inserted{\unsure{discretification}} $\pi'_1$ of $\pi'$ which gives the same oriented prediscretification (resp. the same arithmetization) as $\pi'_0$. Thus we do have $\unsure{P}^{+}(\pi)\longrightarrow P^+(\pi')$, $A(\pi)\longrightarrow A(\pi')$, we still have to express injectivity, which amounts to the relations
\[
\begin{gathered}
(2)\quad
\unsure{\widehat{\mathcal G}^{+}}(\pi_0)
=
\unsure{\widehat{\mathcal G}^{+}}(\pi'_0)\cap
\unsure{\widehat{\widehat{\mathcal G}}}(\pi_0),\qquad
\mathcal M(\pi_0)=\mathcal M(\pi'_0)\cap
\unsure{\widehat{\widehat{\mathcal G}}}(\pi_0)
\end{gathered}
\]
\scanpage{474}
\mspage{240}

i.e. a loop \unsure{automorphism} of $\pi^{+}$ which, on $\pi'$, belongs to $\unsure{\widehat{\mathcal G}}(\pi'_0)$ \deleted{\illegible{2}} resp. to $\mathcal M(\pi'_0)$, is already in $\unsure{\widehat{\mathcal G}}(\pi_0)$ resp. in $\mathcal M(\pi_0)$.

The \deleted{\illegible{1}} marked assertion being admitted, the \inserted{un}marked assertion also means that the canonical homomorphism
\[
(3)\qquad
\mathbb{\Gamma}_a=\mathcal M_a(\pi)/\unsure{\widehat{\mathcal G}}^{+}_a(\pi)
\longrightarrow
\mathbb{\Gamma}_{a'}=\mathcal M_{a'}(\pi')/\unsure{\widehat{\mathcal G}}^{+}_{a'}(\pi')
\]
is \uline{injective}. By construction (via $\mathbb{\Gamma}_{\unsure{\mathbb Q}}$) it is also surjective, so we would find:

\textcircled{2} the canonical homomorphism (3) is bijective

and we would also find a bijection
\[
(4)\qquad P_a(\pi)\longrightarrow P_{a'}(\pi')
\]
between \unsure{arithmetizations} of $\pi$ of type $a$, and arithmetizations of $\pi'$ of type $a'$, \uline{compatible} with the isomorphism \unsure{(3)}.

Finally, a result in the case $\pi'$ characteristically \unsure{invariant} would extend at once to the case of an arbitrary open $\pi'\subset\pi$.

Proceeding \inserted{for example} as in \S 28 by means of ``correspondences'' \inserted{outer arithmetic\unsure{,}} between

\scanpage{475}

``potential curves'', we find a connected groupoid (pointed by the $\widehat{\pi}_{g,\nu}$, for example, \deleted{\unsure{)}} which are \unsure{canonically} isomorphic there), whose outer $\pi_1$ is the common value of the $\mathbb{\Gamma}_{g*}=\mathbb{\Gamma}_{g,\nu}$.

To find \unsure{an} the isomorphism \inserted{\unsure{canonical}} between $\mathbb{\Gamma}_{g,\nu}$ and $\mathbb{\Gamma}_{g',\nu'}$, we choose a \deleted{\unsure{common multiple}} \inserted{correspondence between} a $\pi_{g,\nu}$ and $\pi_{g',\nu'}$, for example by \unsure{regarding} \deleted{\unsure{$N$ of $2g-2+\nu$ and $2g'-2+\nu'$}} both of them (passing if need be: from $g,\nu$ to $g,\widetilde{\nu}$ with $\widetilde{\nu}\geq\nu$, and likewise for $g',\nu'$ to $(g',\widetilde{\nu}')$ with $\widetilde{\nu}'\geq\nu'$) as finite subgroups of the same $\pi_{0,3}$ \unsure{and} whence
\[
\mathbb{\Gamma}_{0,3}\xrightarrow{\sim}\Gamma_{g,\widetilde{\nu}},
\qquad
\Gamma_{0,3}\xrightarrow{\sim}\Gamma_{g',\widetilde{\nu}'}.
\]

We would now like to see what, in the preceding yoga, is independent of any conjecture. For example, the fact that the surjective \unsure{homomorphisms} \inserted{\unsure{canonical}} $\mathbb{\Gamma}_{g,\nu}\unsure{\longrightarrow}\mathbb{\Gamma}_{g,\nu'}$ ($g$ the same, $\nu'$ constant) is not established.

However, if we knew that in the situation of $(\pi,\pi')$ with $\pi'$ characteristically \unsure{invariant} we had
\[
\widehat{\widehat{\mathcal T}}(\pi'_0)\cap
\unsure{\widehat{\mathcal G}}^{+}(\pi_0)
=
\unsure{\widehat{\mathcal G}}^{+}(\pi_0)
\]

\scanpage{476}
\mspage{241}

(which is a statement of a relatively innocuous nature about profinite groups with loops) we would conclude the \uline{bijectivity} of $\mathbb{\Gamma}_{\pi_0}\longrightarrow\mathbb{\Gamma}_{\pi'_0}$ in this situation, and we would conclude from this that the kernel of $\mathbb{\Gamma}_{\unsure{\mathbb Q}}\longrightarrow\mathbb{\Gamma}_{0,3}$ \deleted{\unsure{maps|maps it}} maps \uline{trivially} into the $\mathbb{\Gamma}_{g,\nu}$ for $\forall\nu$ -- i.e. that it maps into the $\mathbb{\Gamma}_{g,\nu}$ through the \uline{intermediary} of $\mathbb{\Gamma}_{0,3}$ \margin{\unsure{$!\! -\! A$}} -- and that moreover $\mathbb{\Gamma}_{0,3}\longrightarrow\Gamma_{g,\nu}$ is
\deleted{\unsure{lack|for lack} of having established this, \unsure{I am going to show} at least that the image of the \unsure{kernels|kernel} $\Theta$ of $\mathbb{\Gamma}_{\unsure{\mathbb Q}}\longrightarrow\mathbb{\Gamma}_{0,3}$ \unsure{in} the $\mathbb{\Gamma}_{g,\nu}$ is \uline{\unsure{finite|finite}} which also means}
\inserted{\uline{an isomorphism} for $\forall(g,\nu)$} with $\nu$ \uline{sufficiently} large -- it suffices that the surface of type $(g,\nu)$ \deleted{\unsure{can}} (or of type $(g,\nu')$ with a $\nu'\leq\nu$) can be obtained as an \unsure{étale} covering of $X_{0,3}$, i.e. that $\pi_{g,\nu'}$ is realized \inserted{(\unsure{with its structure} of loops)} as a \unsure{subgroup|subgroup} \deleted{with loops} \inserted{of finite index} \deleted{\unsure{of in}} of $\pi_{0,3}$. Then, \deleted{\unsure{taking}} \inserted{considering} a $\pi''$ \inserted{$\unsure{\simeq|\sim}\pi_{g'',\nu''}$} $\unsure{\subset|\in}\pi_{g,\nu'}$ of finite index and characteristic in $\pi_{0,3}$, we will have
\scanpage{477}

\diagram{\begin{tikzcd}[ampersand replacement=\&]
\mathbb{\Gamma}_{0,3} \arrow[r,"\sim"] \& \mathbb{\Gamma}_{g'',\nu''} \\
\mathbb{\Gamma}_{\mathbb Q} \arrow[u] \arrow[r] \& \mathbb{\Gamma}_{g,\nu'} \arrow[u,"\unsure{\wr|2}"']
\end{tikzcd}}
commutative \unsure{diagram}.

whence our \unsure{assertion} $\mathbb{\Gamma}_{0,3}\xrightarrow{\sim}\mathbb{\Gamma}_{g,\nu'}$, and a fortiori the surjective homomorphism $\mathbb{\Gamma}_{0,3}\longrightarrow\mathbb{\Gamma}_{g,\nu}$ which \unsure{factors} it \inserted{$(\nu\geq\nu')$} is bijective. Let us note \unsure{for example} that the $\mathbb{\Gamma}_{0,\nu}\longrightarrow\mathbb{\Gamma}_{0,3}$ $(\nu\geq3)$ are then bijective. But \deleted{it \illegible{3}}\inserted{\deleted{\illegible{1}}} note that in any case, we obtain a result for a $g,\nu$ only if $\nu\geq3$ -- so this says nothing about the \unsure{possible} faithfulness \unsure{for example} of $\mathbb{\Gamma}_{0,3}\unsure{\longrightarrow|-}\mathbb{\Gamma}_{g,0}=\mathbb{\Gamma}_g$ $(g\geq2)$ \unsure{or} $\mathbb{\Gamma}_{0,3}\longrightarrow\mathbb{\Gamma}_{1,1}=\mathbb{\Gamma}_1$.

\margin{We no longer make any conjectural hypothesis (\unsure{such} as injectivity in (3)}
Let an algebraic curve \inserted{$\unsure{U}$ of type $g,\nu$} be defined over a subfield \inserted{$K$} finite over $\mathbb Q$ of $\overline{\mathbb Q}_0$, so it defines an outer action of the open subgroup $\Gamma=\Gamma_K\subset\mathbb{\Gamma}_{\mathbb Q}$ on $\pi_1(\unsure{U}_{\overline{\mathbb Q}_0})$. \deleted{Let us suppose} Enlarging $K$ if need be, \deleted{\unsure{to} $K'$, \unsure{we may suppose} \unsure{that}} and \unsure{enlarging}

\scanpage{478}
\mspage{242}

$\nu$ to $\nu'$ (i.e. making holes \inserted{to \unsure{find} $\unsure{U}'$}) \uuline{we may suppose that} $\unsure{U}'$ \uuline{is an étale covering of} \unsure{\deleted{$\mathbb P^1_K$}} $(\unsure{U}_{0,3})_K$.
\margin{we put $\pi'=\pi_1(\unsure{U}'_{\overline{\mathbb Q}_0})$ $\pi=\pi_1(\unsure{U}_{\overline{\mathbb Q}})$}
\deleted{\unsure{Let us choose}} \deleted{\unsure{Let us consid-er} $\pi$ of $\unsure{U}'$ \unsure{relative} to $K$, so $\unsure{\Gamma'}$ operates \unsure{actually} on $\pi$.}
Then let $\unsure{\Theta|\Theta'}$ be the kernel of
\diagram{\begin{tikzcd}[ampersand replacement=\&]
\unsure{\Gamma'|\Gamma} \arrow[r] \arrow[d,phantom,"\cap" description] \& \mathbb{\Gamma}_{0,3} \\
\mathbb{\Gamma}_{\mathbb Q} \&
\end{tikzcd}},
it therefore operates \deleted{by} inner \unsure{automorphisms|automorphisms} on $\widehat{\pi}_{\unsure{03|0,3}}$, so \deleted{\unsure{its}}\inserted{the} image \inserted{of $\Theta'=\Gamma'\unsure{\cap|\wedge}\Theta$} in the group of outer \unsure{automorphisms|automorphisms} of $\pi'\subset\widehat{\pi}_{0,3}$ is a finite subgroup -- a fortiori its image in $\mathbb{\Gamma}_{\pi'}$, and in \inserted{$\mathcal N(\pi)$ and in} $\mathbb{\Gamma}_{\pi}$. The same holds (choosing a \unsure{base|good} point rational over $K$ in $\unsure{U}'$) for the \unsure{actual} operation \unsure{in|of} $\unsure{\pi_1(U,P)|\pi_1(U',P)}$ -- for it suffices to apply the preceding result to $\unsure{U}\unsure{\backslash|-}\{P\}$. We also conclude from this that \deleted{if $(g,\nu)$ is \unsure{such}} the kernels of the homomorphisms $\mathbb{\Gamma}_{0,\nu}\longrightarrow\mathbb{\Gamma}_{0,3}$ $(\nu\geq3)$ are \uline{finite}. We easily see that for every abelian variety $A$ defined over $K$, the operation of $\Theta'$ on the $\prod_\ell T_\ell(A)$ takes place through a finite quotient group.

It \unsure{becomes} very difficult to imagine, how it \unsure{\deleted{$\Theta$}} could happen that $\unsure{\Theta'|\Theta}$ does not operate, in fact, trivially ! \unsure{I have}
\scanpage{479}

even the impression that I can show thanks to the result of the Russian which Deligne pointed out to me, \deleted{\unsure{I think}} that $\mathbb{\Gamma}_{\mathbb Q}\longrightarrow\mathbb{\Gamma}_{0,3}$ is an isomorphism ! \unsure{That|Tha} \unsure{would mean} that the outer actions of $\mathbb{\Gamma}_{\mathbb Q}$ on $\pi_1$'s \inserted{$\simeq\overline{\pi}_{g,\nu}$} can be interpreted (at least for $\nu$ large enough, $g$ being fixed) as splittings \unsure{of} the extension \inserted{$\mathcal N_{g,\nu}$} of $\mathbb{\Gamma}_{g,\nu}$ by $\widehat{\mathcal T}_{g,\nu}$ — or of the extension $\mathcal M_{g,\nu}$ of $\mathbb{\Gamma}_{g,\nu}$ by $\widehat{\mathcal G}_{g,\nu}$, when it is a matter of actual actions.

Of course, the essential question which then arises (admitting $\mathbb{\Gamma}_{\mathbb Q}\simeq\mathbb{\Gamma}_{0,3}$) is to \unsure{characterize} $\mathbb{\Gamma}_{0,3}$ algebraically, as well as the $\mathcal N_{g,\nu}$ — and to give an algebraic description also of the ``admissible'' splittings — i.e. \unsure{corresponding} really and truly to algebraic curves over algebraic number fields.

\scanpage{480}
\mspage{243}

that \deleted{it suffices} these are exactly the actions lifting the \unsure{given|given} outer action, which normalize no loop \unsure{subgroup|subgroups}, \inserted{or again such that $\widehat{\pi}_{0,3}^{\Gamma''+}$ \deleted{\illegible{1}} $(1)$} (But it is \unsure{conceivable} that more subtle conditions may have to be added to it, involving the Frobenius ... ) We may propose to find a description of the \uline{outer} ``admissible'' actions, \deleted{\unsure{in terms}} \inserted{without having} to pass through admissible actual actions — and to \unsure{endea-|ende-}\linebreak\unsure{vor} to give a more or less plausible definition of them. The difficulty of \unsure{course|course} comes from the fact that if $\Gamma''$ operates \uline{by outer automorphisms} on \deleted{\illegible{1}} $\widehat{\pi}_{0,3}$, it does not thereby operate by outer automorphisms on a given open subgroup \inserted{$\pi'$} of $\widehat{\pi}_{0,3}$. But we are going to \unsure{suppose} precisely that such an action \uline{exists}, \deleted{\unsure{compatible}} \inserted{in such a way that} \deleted{\unsure{with}} $\pi'\longrightarrow\widehat{\pi}_{0,3}$ is compatible with the outer actions, and that the action of $\Gamma''$ on $\pi$ is \unsure{deduced} (modulo isomorphism) from
